% =====================================================================
%  Guía de Laboratorio — Institución Universitaria de Barranquilla (IUB)
%  Plantilla en LaTeX clonada de "Plantilla guias de laboratorio.docx"
%
%  Compilar con:  pdflatex -> biber -> pdflatex -> pdflatex
%  (En Overleaf: compilador pdfLaTeX; biber se ejecuta solo.)
%
%  Estructura de archivos:
%    guia-laboratorio.tex
%    guialabiub.sty
%    referencias.bib
%    img/encabezado.png   img/pie-pagina.png   img/marca-agua.jpg
% =====================================================================
\documentclass[12pt]{article}

% Opciones: [sinmarca] quita la marca de agua; [borrador] muestra márgenes
%           [ieee] referencias IEEE (por defecto); [apa] referencias APA 7
\usepackage{guialabiub}

% --------------------------- Bibliografía ----------------------------
\addbibresource{referencias.bib}

% --------------------------- Datos de la guía -------------------------
\facultad{}
\modulo{}
\docente{}

\begin{document}

\datosguia
\lineablanco

% ---------------------------------------------------------------------
\section{Objetivos}
Describa los objetivos generales de la práctica de laboratorio. ¿Qué
habilidades o conceptos se espera que los estudiantes aprendan o apliquen
durante la práctica?

% ---------------------------------------------------------------------
\section{Materiales y Equipos}
Enumere todos los materiales y equipos necesarios para llevar a cabo la
práctica de laboratorio. Asegúrese de incluir cualquier equipo
especializado, reactivos químicos, instrumentos de medición, etc.

% ---------------------------------------------------------------------
\section{Procedimiento}
Detalle paso a paso el procedimiento que los estudiantes deben seguir
durante la práctica de laboratorio. Es importante que cada paso sea claro
y conciso, indicando cualquier medida de seguridad relevante y los tiempos
aproximados necesarios para cada etapa.
\lineablanco
\begin{pasos}
  \item Paso 1:
  \item Paso 2:
  \item Paso 3:
\end{pasos}

% ---------------------------------------------------------------------
\section{Cuestiones a Resolver}
Incluya preguntas o problemas que los estudiantes deben resolver como
parte de la práctica de laboratorio. Estas preguntas pueden estar
relacionadas con los resultados esperados, la interpretación de datos, el
análisis de errores, etc.

% ---------------------------------------------------------------------
\section{Análisis de Resultados}
Instrucciones sobre cómo los estudiantes deben analizar y presentar los
resultados obtenidos durante la práctica de laboratorio. Esto puede
incluir la preparación de tablas, gráficos, interpretación de datos y
conclusiones.

% ---------------------------------------------------------------------
\section{Conclusiones}
Pida a los estudiantes que resuman las conclusiones principales de la
práctica de laboratorio. ¿Qué aprendieron? ¿Hubo algún resultado
inesperado?

% ---------------------------------------------------------------------
\section{Recomendaciones de Seguridad}
Recuerde a los estudiantes cualquier medida de seguridad relevante que
deban seguir durante la práctica de laboratorio. Esto puede incluir el uso
de equipos de protección personal, manejo seguro de reactivos químicos,
etc.

% ---------------------------------------------------------------------
\section{Bibliografía}
Incluya una lista de referencias o fuentes consultadas para la
preparación de la guía de laboratorio. Esto puede incluir libros de texto,
artículos científicos o cualquier otra fuente relevante.

% Cite las fuentes en el texto con \cite{clave}. \nocite{*} incluye todas
% las entradas del .bib aunque no se citen; elimínelo si solo quiere las citadas.
\nocite{*}
\bibliografia

\end{document}
